\documentclass[11pt,a4paper]{article}
\usepackage[utf8]{inputenc}
\usepackage[margin=1in]{geometry}
\usepackage{amsmath,amssymb,amsfonts,amsthm}
\usepackage{bm}
\usepackage{booktabs}
\usepackage{tabularx}
\usepackage{graphicx}
\usepackage{hyperref}
\usepackage{cite}
\usepackage{microtype}
\usepackage{algorithm}
\usepackage{algpseudocode}

\hypersetup{
    colorlinks=true,
    linkcolor=blue,
    citecolor=red,
    urlcolor=blue
}

\newtheorem{theorem}{Theorem}
\newtheorem{lemma}{Lemma}
\newtheorem{definition}{Definition}
\newtheorem{remark}{Remark}

\title{\textbf{Deep Early-Exercise Premium Physics-Informed Neural Networks (Deep-EEP-PINN): High-Dimensional American Basket Option Free-Boundary Valuation up to 50 Dimensions}}

\author{
    \textbf{Kartikey Singh}\thanks{Department of Mathematics, University of Delhi, Delhi, India. Email: \texttt{kartikeysingh525@protonmail.com}}
}

\date{\today}

\begin{document}

\maketitle

\begin{abstract}
Valuing high-dimensional American basket options represents a challenging parabolic free-boundary obstacle partial differential equation (PDE) problem. Conventional grid-based solvers (e.g., Finite Difference, Projected Successive Over-Relaxation) suffer from Bellman's curse of dimensionality, while standard Physics-Informed Neural Networks (PINNs) fail due to the non-differentiable payoff kink singularity at expiration $t=T$ and Dirac-delta blow-up of spatial Gamma ($\Gamma$). In this paper, we introduce \textbf{Deep-EEP-PINN}, a mesh-free scientific machine learning architecture for multi-asset American option pricing up to $d=50$ correlated assets ($1,225$ pairwise correlation terms). Deep-EEP-PINN combines the analytical Kim (1990) / Carr-Jarrow-Myneni (1992) Early Exercise Premium (EEP) decomposition with a hard-constraint boundary-encoded neural ansatz that identically satisfies $e_\theta(\mathbf{S}, T) \equiv 0$, eliminating the terminal kink singularity by construction. For multi-asset cross-diffusion, we introduce an $\mathcal{O}(d)$ directional autograd Hessian trace contraction combined with chunked loss backpropagation, constraining peak GPU VRAM to $< 3\,\text{GB}$ on standard accelerators. Rigorous empirical validation against 100,000-path Longstaff-Schwartz Least Squares Monte Carlo (LSM) demonstrates that Deep-EEP-PINN achieves mean pricing errors of \text{Rs. }0.06 (6 paise) in 30D (Dow Jones scale) and \text{Rs. }0.08 (8 paise) in 50D (Nifty 50 scale) with an inference latency of $3.94\,\text{ms}$, delivering a $\mathbf{176,760\times}$ speedup over 100k-path LSM regressions.
\end{abstract}

\noindent\textbf{Keywords:} Physics-Informed Neural Networks (PINNs), Free-Boundary PDEs, American Basket Options, Optimal Stopping, High-Dimensional Scientific Machine Learning, Curse of Dimensionality.

\section{Introduction}
American-style financial derivatives permit the holder to exercise the contract early at any stopping time $\tau \in [0, T]$ prior to maturity $T$. The continuous-time valuation of multi-asset American basket options under correlated Geometric Brownian Motion constitutes a parabolic obstacle problem governed by a Linear Complementarity Problem (LCP):
\begin{equation}
\label{eq:lcp}
\min\left( -\frac{\partial V}{\partial t} - \mathcal{L}_{\text{BS}} V, \; V(\mathbf{S}, t) - h(\mathbf{S}) \right) = 0, \quad (\mathbf{S}, t) \in \mathbb{R}_+^d \times [0, T),
\end{equation}
subject to the terminal obstacle condition $V(\mathbf{S}, T) = h(\mathbf{S})$, where $h(\mathbf{S}) = \max\left( K - \phi(\mathbf{S}), 0 \right)$ is the intrinsic payoff, $\phi(\mathbf{S})$ is the basket aggregation function, and $\mathcal{L}_{\text{BS}}$ denotes the multi-asset Black-Scholes spatial diffusion operator:
\begin{equation}
\mathcal{L}_{\text{BS}} V = \frac{1}{2} \sum_{i=1}^d \sum_{j=1}^d \rho_{ij} \sigma_i \sigma_j S_i S_j \frac{\partial^2 V}{\partial S_i \partial S_j} + r \sum_{i=1}^d S_i \frac{\partial V}{\partial S_i} - r V.
\end{equation}

In one spatial dimension ($d=1$), classical numerical methods such as Projected Successive Over-Relaxation (PSOR) and Crank-Nicolson finite difference schemes achieve high precision \cite{cryer1971solution, brennan1977valuation}. However, as the asset dimension increases ($d \ge 5$), grid-based discretization collapses due to the curse of dimensionality, requiring $\mathcal{O}(N^d)$ spatial nodes. 

While recent deep learning approaches based on Backward Stochastic Differential Equations (BSDEs) and deep optimal stopping \cite{han2018solving, becker2019deep, becker2020solving} offer path-wise approximations, they remain simulation-dependent, require thousands of Monte Carlo samples per evaluation point, and do not provide continuous closed-form spatial derivatives (Greeks) via automatic differentiation. 

Conversely, standard Physics-Informed Neural Networks (PINNs) \cite{raissi2019physics, al2018solving} suffer from severe optimization stalling and Dirac-delta Gamma blow-up at expiration due to the non-smooth kink in $\max(K - S, 0)$.

\section{Methodology: The Deep-EEP-PINN Architecture}

\subsection{Analytical Early Exercise Premium Decomposition}
Following the analytical foundations established by Kim \cite{kim1990analytic} and Carr, Jarrow, and Myneni \cite{carr1992alternative}, the total American option value is decomposed into an exact European anchor and an early exercise premium:
\begin{equation}
V_{\text{Amer}}(\mathbf{S}, t) = V_{\text{Euro}}^{\text{exact}}(\mathbf{S}, t) + e_\theta(\mathbf{S}, t),
\end{equation}
where $V_{\text{Euro}}^{\text{exact}}(\mathbf{S}, t)$ is evaluated via exact closed-form Black-Scholes formulas (or lognormal moment-matched approximations for arithmetic baskets \cite{milevsky1998valuing, gentle1993basket}).

\subsection{Hard-Constraint Boundary-Encoded Neural Ansatz}
To guarantee $e_\theta(\mathbf{S}, T) \equiv 0$ strictly without training error, the premium is parameterized as:
\begin{equation}
e_\theta(\mathbf{S}, t) = \text{Softplus}\left(\mathcal{N}_\theta(\mathbf{S}, t)\right) \cdot \left(\frac{T-t}{T}\right) \cdot \prod_{i=1}^d \left(1 - \frac{S_i}{S_{\max}}\right) \cdot K\left(1 - e^{-rT}\right).
\end{equation}

\subsection{$\mathcal{O}(d)$ Directional Autograd Trace Contraction}
To evaluate the cross-diffusion term $\frac{1}{2} \operatorname{Tr}\left( \boldsymbol{\Sigma}(\mathbf{S}) \mathbf{H} \right)$ without constructing the $d \times d$ Hessian matrix, we decompose the covariance matrix into its Cholesky factor $\mathbf{L}$ ($\boldsymbol{\Sigma} = \mathbf{L} \mathbf{L}^T$) and compute directional second derivatives:
\begin{equation}
\frac{1}{2} \sum_{k=1}^d \mathbf{v}_k^T \nabla^2_{\mathbf{S}} e_\theta \, \mathbf{v}_k, \quad \mathbf{v}_k = \mathbf{L}_k \odot \mathbf{S}.
\end{equation}

\section{Empirical Results \& Benchmark Suite}

All experiments were executed on an NVIDIA Tesla P100 GPU (16GB VRAM) and validated against 100,000-path Longstaff-Schwartz Least Squares Monte Carlo regressions.

\begin{table}[htbp]
\centering
\small
\caption{Master Benchmark Results across All Dimensions ($d=1, 5, 10, 30, 50$).}
\begin{tabularx}{\textwidth}{@{}l c c c c c@{}}
\toprule
\textbf{Problem Setup} & \textbf{Ground Truth} & \textbf{EEP-PINN Error} & \textbf{Inference Latency} & \textbf{Speedup vs Truth} \\ \midrule
1D American Put ($d=1$) & PSOR (90k nodes) & $L_2 = 0.37\%$, MAE = \text{Rs. }0.08 & $11.84\,\text{ms}$ & $97.1\times$ vs PSOR \\
5D Geometric Basket ($d=5$) & 100k LSM ($\pm 1.96\text{SE}$) & Diff = \text{Rs. }0.06 (ATM) & $3.50\,\text{ms}$ & $2712.2\times$ vs LSM \\
5D Arithmetic Basket ($d=5$) & 100k LSM ($\pm 1.96\text{SE}$) & Diff = \text{Rs. }0.02 (ATM) & $7.55\,\text{ms}$ & $1010.0\times$ vs LSM \\
10D High-Dim Basket ($d=10$) & 100k LSM ($\pm 1.96\text{SE}$) & Diff = \text{Rs. }0.14 (Mean) & $3.47\,\text{ms}$ & $6882.6\times$ vs LSM \\
30D Dow Jones Scale ($d=30$) & 100k LSM (498 basis) & Diff = \text{Rs. }0.06 (Mean) & $3.26\,\text{ms}$ & $62,842.1\times$ vs LSM \\
50D Nifty 50 Scale ($d=50$) & 100k LSM (1328 basis) & Diff = \text{Rs. }0.08 (Mean) & $3.94\,\text{ms}$ & $\mathbf{176,760.1\times}$ vs LSM \\ \bottomrule
\end{tabularx}
\end{table}

\section{Conclusion}
Deep-EEP-PINN establishes a scalable, mesh-free framework for continuous-time free-boundary obstacle PDEs in mathematical finance, breaking the curse of dimensionality up to 50 correlated assets while providing real-time millisecond inference for high-frequency risk management.

\begin{thebibliography}{99}
\bibitem{kim1990analytic} I. J. Kim, \emph{The Analytic Valuation of American Options}, The Review of Financial Studies, 3(4):547--572, 1990.
\bibitem{carr1992alternative} P. Carr, R. Jarrow, and R. Myneni, \emph{Alternative Methods for Valuing American Options}, Finance and Stochastics, 2(1):87--106, 1992.
\bibitem{raissi2019physics} M. Raissi, P. Perdikaris, and G. E. Karniadakis, \emph{Physics-informed neural networks}, Journal of Computational Physics, 378:686--707, 2019.
\bibitem{longstaff2001valuing} F. A. Longstaff and E. S. Schwartz, \emph{Valuing American Options by Simulation: A Simple Least-Squares Approach}, The Review of Financial Studies, 14(1):113--147, 2001.
\bibitem{han2018solving} J. Han, A. Jentzen, and W. E, \emph{Solving high-dimensional partial differential equations using deep learning}, PNAS, 115(34):8505--8510, 2018.
\bibitem{becker2019deep} S. Becker, P. Cheridito, and A. Jentzen, \emph{Deep Optimal Stopping}, Journal of Machine Learning Research, 20(74):1--25, 2019.
\bibitem{becker2020solving} S. Becker, P. Cheridito, A. Jentzen, and T. Welti, \emph{Solving High-Dimensional Optimal Stopping Problems Using Deep Learning}, SIAM J. Sci. Comput., 42(1):B1--B29, 2020.
\bibitem{cryer1971solution} C. W. Cryer, \emph{The Solution of a Quadratic Programming Problem using Systematic Overrelaxation}, SIAM J. Control, 9(3):385--392, 1971.
\bibitem{brennan1977valuation} M. J. Brennan and E. S. Schwartz, \emph{The Valuation of American Put Options}, The Journal of Finance, 32(2):449--462, 1977.
\bibitem{milevsky1998valuing} M. A. Milevsky and S. E. Posner, \emph{Valuing Exotic Options by Approximating the Swapping Density}, JFQA, 33(3):409--425, 1998.
\bibitem{gentle1993basket} D. Gentle, \emph{Basket Options}, Risk, 6(12):51--52, 1993.
\bibitem{al2018solving} A. Al-Aradi, A. Correia, D. de Freitas, and P. Garnier, \emph{Solving nonlinear and high-dimensional PDEs in finance using PINNs}, arXiv:1811.08782, 2018.
\end{thebibliography}

\end{document}
